\section{Организация исполнения и системная оценка энергоэффективности}

\subsection{Разреженность и условное выполнение}

Сокращение объёма работы при инференсе может основываться
на структуре параметров и на свойствах текущего входа.
Разреженное представление позволяет исключать действия
над отсутствующими элементами. Условное выполнение изменяет
состав обрабатываемых частей модели по результатам
промежуточного вычисления. Эти подходы различаются
по условиям корректности и способу управления исполнением.

В ускорителе Efficient Inference Engine (EIE), представленном С.~Ханем
и соавторами в 2016 году, вычисление выполнялось непосредственно
над сжатой моделью с разреженными весами и совместным
использованием значений параметров~\cite{han2016eie}.
Аппаратура обрабатывала индексы ненулевых элементов и пропускала нулевые
активации. Сокращение объёма хранения сочеталось с возможностью выполнять
вычисления непосредственно над сжатыми данными.

Эффект разреженности зависит от затрат на индексы,
передачу операндов и распределение работы между узлами.
Число ненулевых элементов характеризует алгоритмическую
структуру, но не задаёт загрузку аппаратуры. Для оценки
энергии требуется учитывать операции с данными и метаданными,
а также действия управляющих механизмов.

В BranchyNet С.~Тирапиттаянон, Б.~Макдэнел и Х.~Т.~Кунг
в 2016 году исследовали раннее завершение обработки
в классификационных сетях~\cite{teerapittayanon2016branchynet}.
Дополнительные выходные ветви позволяли получать ответ
до прохождения всех слоёв, если выполнялся установленный
критерий уверенности. Состав операций зависел от входного
примера; для остальных входов сохранялась последующая обработка.

Ранний выход требует совместного выбора критерия завершения
и допустимого изменения качества. Его затраты включают
вычисление промежуточных ответов и проверку условия.
Среднее сокращение объёма обработки зависит от состава
входных данных. Поэтому выводы о сокращении работы при условном исполнении зависят от
распределения входных данных и установленных требований к ответу.

В Bit-Pragmatic 2016 года использовался пропуск нулевых
битовых слагаемых операнда~\cite{pragmatic2016}.
При таком представлении операндов система исключала операции с нулевым вкладом
в результат. Этот механизм отличается
от приближённого изменения модели или раннего завершения:
в первом случае исключается нулевой вклад, а во втором требуется проверить
допустимость приближённого ответа.

\subsection{Согласование параллелизма, обмена данными и управления}

Энергетическая эффективность исполнительной системы зависит
от согласования арифметических узлов, памяти и управления.
Параллелизм сокращает время выполнения при достаточной
подаче операндов и распределении работы. Неравномерность
нагрузки, преобразование представлений и синхронизация
создают дополнительные затраты. Эти зависимости относятся
к различным видам вычислительных архитектур.

Изменение разрядности позволяет учитывать численные свойства обрабатываемых
данных при организации работы аппаратуры.
В Stripes П.~Джадда и соавторов 2016 года использовалась
последовательная обработка битов с предварительно выбранной
для слоя точностью~\cite{stripes2016}. В Dynamic Stripes
А.~Дельмас и соавторы в 2017 году дополнили этот подход
определением значащих разрядов текущих групп активаций~\cite{dynamic2017}.
В отличие от заранее выбранной точности слоя, здесь решение принималось для
текущей группы активаций по их значениям.

В Bit Fusion Х.~Шармы и соавторов 2018 года вычислительные
узлы составлялись из меньших битовых элементов~\cite{bitfusion2018}.
Архитектура была представлена на Международном симпозиуме
по архитектуре компьютеров (International Symposium on
Computer Architecture, ISCA). Состав объединяемых элементов
зависел от разрядности операции. Перестройка исполнительных
ресурсов позволяла изменять структуру параллельной обработки.

Эти проекты демонстрируют разные способы изменения работы аппаратуры в
зависимости от разрядности данных. Организация обмена и выбор формата данных
согласуются также с ёмкостью буферов и порядком выполнения
операций. Если преобразование требует дополнительных
проходов или хранения метаданных, их стоимость включается
в оценку полного исполнения. Сокращение арифметических
действий сопоставляется с изменением остальных затрат.

Выбор размера группы данных влияет на точность настройки режима и затраты
управления. Если один режим задаётся для большой группы, различия отдельных
операндов учитываются хуже. При уменьшении группы решения приходится принимать
чаще, что увеличивает объём управляющих действий. При проектировании разреженных,
низкоразрядных и условно исполняемых систем размер группы
выбирается с учётом этих затрат и устройства аппаратуры.

\subsection{Визуальные трансформеры и диффузионные модели}

В начале 2020-х годов развивались методы обработки изображений
на основе трансформерных архитектур. В Vision Transformer,
описанном А.~Досовицким и соавторами в работе 2020 года
и представленном на конференции ICLR в 2021 году, изображение обрабатывалось
как последовательность фрагментов~\cite{vit2021}.
Это изменило состав операций и характер взаимодействия элементов
представления. Наличие матричных умножений сохранялось,
но наличие одинакового типа операций не означало, что требования к исполнению
совпадали с требованиями свёрточной сети.

В диффузионных вероятностных моделях Дж.~Хо, А.~Джайна
и П.~Аббила 2020 года синтез строился через последовательность
шагов удаления шума~\cite{ho2020}. Формирование одного
изображения включало повторное выполнение обученной модели.
Повторное применение модели необходимо учитывать при оценке численной
точности. Последствия изменения арифметики проверяются для всего процесса
генерации, поскольку результат каждого шага используется на следующем.

Работа У.~Пиблса и С.~Се 2022--2023 годов о диффузионных
трансформерах соединила пошаговое удаление шума с трансформерной архитектурой~\cite{dit2023}. При оценке энергоэффективности необходимо учитывать эти изменения: меняются размеры
промежуточных представлений, число шагов и состав операций.
Применимость результатов сокращения точности, полученных
на классификационных сетях, к генеративным моделям требует
отдельного экспериментального обоснования.

Появление новых классов нейросетевых моделей изменяло требования к аппаратным
средствам их исполнения.
Матричные операции сохранялись в составе вычислений,
но изменялись их размеры, взаимосвязи и порядок исполнения.
Разработанные ранее аппаратные подходы могли применяться к новым моделям при
условии, что организация обмена данными и свойства арифметики соответствовали
составу и порядку выполняемых операций.

\subsection{Сценарии применения и организация нагрузки}

Методы энергоэффективного исполнения разрабатывались с учётом разных условий
применения нейронных сетей. Встраиваемое
устройство, интерактивный сервис и обработка заранее
подготовленного набора данных предъявляют различные требования
к памяти, времени ответа и организации параллелизма.
Даже одна и та же нейронная сеть может использоваться в этих условиях
по-разному.
Выбор способа исполнения включает определение порядка
поступления запросов и допустимого ожидания результата.

Мобильные и встраиваемые приложения стали одной из областей
развития компактных архитектур. В MobileNets 2017 года
ресурсные ограничения учитывались посредством изменения
структуры сети и её параметров~\cite{howard2017mobilenets}.
В этом случае часть решений, влияющих на затраты инференса, принимается уже
при проектировании модели.
При последующем исполнении сохраняется необходимость
согласования операций с доступной памятью и библиотеками.

Для серверных приложений существенны ограничения
на время ответа отдельных запросов. В анализе тензорного
процессора Google 2017 года условия сравнения включали
требования к 99-му процентилю задержки~\cite{jouppi2017}.
Такой показатель характеризует границу, в пределах которой
должна завершаться обработка заданной доли запросов.
Повышение средней пропускной способности само по себе не гарантирует
соблюдения этого ограничения. Следовательно, организация
исполнения оценивается с учётом распределения времени ответа.

Различия сценариев были формализованы в MLPerf Inference
2019--2020 годов~\cite{mlperf2020}. Методика разделяла
последовательное поступление отдельных запросов,
обработку нескольких потоков, серверную нагрузку
и обработку заранее доступного набора данных.
Таким образом, при описании измерения указывался также способ поступления
запросов. Сравнение устройств предполагало
соответствие не только модели, но и сценарию её использования.

При измерении энергии также необходимо указать, как поступают запросы и
сколько из них обработано за выбранный интервал. Увеличение пакета может улучшать
повторное использование данных, но добавлять ожидание
формирования пакета. При нерегулярном поступлении запросов
в измеряемый интервал могут входить периоды неполной загрузки.
Поэтому энергия выполнения установленного числа запросов
сопоставляется при определённых требованиях к их обслуживанию.
Следовательно, снижать энергозатраты можно также за счёт организации обработки
запросов, наряду с изменением модели, численного формата и аппаратуры.

\subsection{Критерии сравнения и границы исторических выводов}

Рост числа ускорителей потребовал общих правил сравнения.
В публикации о MLPerf Inference, представленной на ISCA
в 2020 году, разработчики описали методику воспроизводимого
сопоставления различных систем инференса~\cite{mlperf2020}.
Для сравнения задаются выбранная модель, требования
к качеству, сценарий поступления запросов и способ измерения
производительности. Если эти условия различаются, сравниваемые показатели могут характеризовать
разные вычислительные задачи.

В публикациях об энергоэффективной аппаратуре сопоставлялись
затраты арифметического узла, кристалла и всей системы.
Эти показатели характеризовали разные объекты исследования.
По мере усложнения устройств возрастало значение передачи
данных и работы обслуживающих компонентов.
Сравнение отдельных операций дополнялось оценкой выполнения
прикладной задачи в выбранной системе.

Различие мощности и энергии также влияло на содержание
публикуемых результатов. Сокращение времени ответа,
увеличение пропускной способности и снижение энергии
не означали одного и того же достижения.
В работах об ускорителях производительность рассматривалась
совместно с затратами исполнения и требованиями приложения.
Такое сопоставление позволяло обсуждать пользу специализации
в конкретных условиях применения.

В численных исследованиях сопоставлялись разные характеристики
результата: совпадение вычисленных значений, погрешности
отдельных слоёв и качество решения прикладной задачи.
Эти характеристики отвечали на разные исследовательские
вопросы. Расширение классов приложений также меняло
содержание сравнения: для генеративных моделей
рассматривалось качество распределения результатов,
а не только доля правильных ответов в классификации.

Одним из способов связать энергетические затраты
с выполненной работой стала оценка энергии на обработанный
запрос. Её смысл зависел от требований к качеству и задержке,
сценария нагрузки и учитываемых частей системы.
Развитие таких критериев сопровождало переход от описания
возможностей отдельного блока к сравнению прикладных
систем инференса.

Число параметров и операций описывает модель, а поддержка
низкоразрядной арифметики — возможности аппаратуры.
Фактическая энергия дополнительно зависит от передачи данных,
преобразования представлений и управления. Их учёт связывает
изменение модели с характеристиками полного приложения.

\subsection{Взаимосвязь алгоритмических и аппаратных подходов}

Методы повышения энергоэффективности изменяют разные части вычислительного
процесса. Компактная архитектура изменяет
структуру модели; квантование — представление значений;
организация памяти и расписания — способ выполнения операций.
Условное исполнение дополнительно связывает объём обработки
с текущими входами. Поэтому при сравнении методов необходимо указать, какие именно затраты
меняются при их применении.

На алгоритмическом уровне заданное качество может достигаться
моделью с меньшим объёмом параметров или операций.
Дистилляция и проектирование компактных сетей позволяют получать такие модели
разными способами. Однако уменьшение размера модели само по себе не доказывает
снижения энергии. Для этого необходимо оценить фактические действия системы
при выбранном способе исполнения.

На аппаратном уровне размещение повторно используемых
данных уменьшает потребность в передаче между уровнями
памяти. Энергетический эффект зависит от ёмкости хранения,
структуры доступа и свойств нагрузки. В Eyeriss эта связь
исследовалась для потоков данных свёрточных вычислений~\cite{eyeriss2016}.
Регулярные массивы, разреженные исполнительные структуры
и программируемые процессоры различаются по способам повторного использования данных и организации доступа
к памяти.

Совместное применение методов требует оценки их взаимодействия.
После сокращения арифметических действий ограничивающим
фактором может стать передача данных. Дополнительные коды
или промежуточные преобразования могут увеличивать объём
обработки при уменьшении исходного представления.
Поэтому суммарную экономию нельзя получить простым сложением результатов
отдельных преобразований без проверки их взаимодействия.

Критерий энергетической эффективности относится к результату
полного исполнения при установленных требованиях к качеству
и времени ответа. При этих условиях можно сравнивать разные классы методов, применяемых в
различных вычислительных системах. По мере развития энергоэффективных
вычислений исследователи учитывали больше видов затрат и совместно
рассматривали выбор алгоритма, формата данных и организации аппаратуры.

\subsection{Воспроизводимость экспериментов}

Для сравнения ускорителей необходимо описать условия, в которых были получены
их характеристики. Воспроизведение опыта требует
описания подготовки параметров модели, представления входных данных
и процедуры вычисления выходного ответа. Для адаптивных систем
дополнительно определяется правило выбора режима. Такое описание позволяет проверить, соответствует ли повторённый эксперимент
опубликованному методу.

Историческим примером ответа на эту проблему служит разработка
MLPerf Inference. В публикации 2020 года сравнение систем
связывается с общими правилами, сценариями нагрузки
и проверкой качества~\cite{mlperf2020}. Общие правила позволяют сравнивать разные реализации при заданных условиях
нагрузки и качества. Устройства при этом сохраняют свои архитектурные
различия, а измеряемые показатели характеризуют выполнение одной и той же
установленной задачи.

Воспроизведение эксперимента позволяет проверить процедуру
получения результата и его объяснение. Изменение отдельной
предпосылки, например размера пакета или способа выбора
масштаба, используется для оценки устойчивости вывода
к соответствующему фактору. Чтобы объяснить наблюдаемый эффект работой аппаратуры, необходимо учитывать
все условия эксперимента. Совпадение результатов повторных запусков при
неизменных настройках ещё не показывает, сохранится ли вывод при изменении
этих условий.

\begin{samepage}
Измерения на изготовленном устройстве, расчёт проектируемой схемы и
программный эксперимент описывают разные этапы исследования. При сравнении
опубликованных показателей проверяется, совпадают ли измеряемые величины,
условия исполнения и учитываемые части системы. Для каждого числа необходимо
знать способ его получения и принятые допущения.
\par\nopagebreak[4]

Развитие методов оценки требовало более подробного описания экспериментов.
Помимо устройства архитектуры указывались режимы работы, ограничения и условия
сравнения. Эти сведения позволяют различать повторную проверку опубликованного
результата, его проверку в новых условиях и испытание другого способа
вычислений. Тем самым уточнялись требования к экспериментам, подтверждающим
снижение энергозатрат.
\par
\end{samepage}
